%% agent_9_fm_configurations_e3.tex
%% Chain-level Fulton--MacPherson configuration-space model for the E_3 operad,
%% as it enters the Stage-1 PhiFA_3 construction in Vol III.
%%
%% Voices: Kapranov, Kontsevich, Soibelman.
%% Standalone working notes; not included in main.tex.

\documentclass[11pt]{article}
\usepackage[margin=1in]{geometry}
\usepackage{amsmath, amssymb, amsthm, mathtools}
\usepackage{enumitem}

\newtheorem{theorem}{Theorem}[section]
\newtheorem{proposition}[theorem]{Proposition}
\newtheorem{lemma}[theorem]{Lemma}
\newtheorem{corollary}[theorem]{Corollary}
\theoremstyle{definition}
\newtheorem{definition}[theorem]{Definition}
\newtheorem{example}[theorem]{Example}
\newtheorem{construction}[theorem]{Construction}
\theoremstyle{remark}
\newtheorem{remark}[theorem]{Remark}
\newtheorem{attack}[theorem]{Attack}
\newtheorem{heal}[theorem]{Heal}
\newtheorem{survivor}[theorem]{Surviving core}

\DeclareMathOperator{\Conf}{Conf}
\DeclareMathOperator{\Aut}{Aut}
\DeclareMathOperator{\Ran}{Ran}
\DeclareMathOperator{\Op}{Op}
\DeclareMathOperator{\End}{End}
\DeclareMathOperator{\Hom}{Hom}
\DeclareMathOperator{\Fact}{Fact}
\DeclareMathOperator{\HH}{HH}
\DeclareMathOperator{\CH}{CH}
\DeclareMathOperator{\Gerst}{Gerst}
\DeclareMathOperator{\Graph}{Graph}
\DeclareMathOperator{\GC}{GC}
\DeclareMathOperator{\GRT}{GRT}

\newcommand{\FMbar}{\overline{F}}
\newcommand{\FM}{\mathcal{FM}}
\newcommand{\cM}{\mathcal{M}}
\newcommand{\Ddisk}{D}
\newcommand{\Ainf}{A_{\infty}}
\newcommand{\Einf}{E_\infty}
\newcommand{\Eone}{E_1}
\newcommand{\Etwo}{E_2}
\newcommand{\PhiFA}{\Phi^{\mathrm{FA}}}
\newcommand{\SpCh}{\mathrm{Sp}^{\mathrm{ch}}}
\newcommand{\CY}{\mathrm{CY}}

\title{Chain-level Fulton--MacPherson configuration-space model for the \texorpdfstring{$E_3$}{E-3} operad}
\author{Agent 9 --- wave \texttt{cfg2026}}
\date{}

\begin{document}
\maketitle

\begin{abstract}
The Stage-$1$ functor $\PhiFA_3$ of the $\CY$-to-chiral correspondence
$\Phi_3 = \SpCh_{\Sigma_2, C} \circ \PhiFA_3$ takes a cyclic $\Ainf$
Calabi--Yau category of dimension three and produces an $E_3$-holomorphic
factorisation algebra on the ambient $\CY_3$ variety. The $E_3$ operad
appearing on the output side is the chain-level Fulton--MacPherson
operad $\FM_3 = \{\FMbar(k, \mathbb{R}^3)^{\FM}\}_{k \geq 1}$: a tower
of smooth manifolds-with-corners whose interior is the configuration
space $\Conf_k(\mathbb{R}^3)$, compactified by iterated real blow-ups
along pairwise and higher-order diagonals. Kontsevich $E_3$-formality
identifies the chain operad $C_*(\FM_3; \mathbb{Q})$ with the
Gerstenhaber$_3$ operad of shifted-Poisson algebras up to a
$\GRT_1(\mathbb{Q})$-torsor. Willwacher's theorem writes this torsor
as $H^0$ of the Kontsevich graph complex $\GC_3$. These notes record
the chain-level operadic content explicitly: boundary strata, the
Boardman--Vogt $W$-construction, the graph-operad comparison, and
the compatibility with the two-stage factorisation of the
correspondence programme.
\end{abstract}

\tableofcontents

%==============================================================================
\section{Cycle 1: ATTACK --- the $E_3$ operad as black box}
%==============================================================================

\begin{attack}[$E_3$ as a symbol without content]
\label{att:cycle-1}
In the Vol III chapter \texttt{cy\_categories.tex} the phrase
``$E_3$-holomorphic factorisation algebra'' appears as the output
target of Stage-$1$ $\PhiFA_3$ (Theorem~\texttt{phi-two-stage-factorisation}).
The $E_3$ operad is named; its operations are not exhibited. In
particular, the following questions have no answer in the chapter as
presently written:
\begin{itemize}
 \item What are the \emph{spaces} $E_3(k)$ of arity-$k$ operations?
 \item What are their \emph{boundary strata}?
 \item What is the chain-level model $C_\bullet(E_3(k); \mathbb{Q})$
       that Kontsevich--Tamarkin formality acts on?
 \item How do the spaces $E_3(k)$ compose operadically?
 \item How does the Stage-$1$ functor $\PhiFA_3$ \emph{use} these
       operations to assemble local observables on a $\CY_3$?
\end{itemize}
Without concrete answers, ``$E_3$-holomorphic factorisation algebra''
is a label, not a mathematical object.
\end{attack}

\begin{survivor}[\ref{att:cycle-1}]
The attack is sharp: the Chriss--Ginzburg standard forbids labels
without content. The operadic data on the output side of $\PhiFA_3$
must be made explicit at the chain level. The Fulton--MacPherson
compactification $\FMbar(k, \mathbb{R}^3)^{\FM}$
\cite[Ann.\ of Math.\ 139, 1994, \S 1]{FultonMacPherson1994}
supplies a canonical model admitting all four pieces of structure:
spaces, boundary strata, chains, operadic composition.
\end{survivor}

%==============================================================================
\section{Cycle 2: HEAL --- configuration spaces and their compactification}
%==============================================================================

\begin{definition}[Ordered configuration space]
\label{def:conf-space}
Let $M$ be a smooth manifold of dimension $n$ and let $k \geq 1$.
The \emph{ordered configuration space} of $k$ points in $M$ is
\[
 \Conf_k(M) \;=\; F(k, M) \;=\; \bigl\{(x_1, \ldots, x_k) \in M^k :
  x_i \neq x_j \text{ for } i \neq j \bigr\}.
\]
It is a smooth manifold of dimension $kn$, open in $M^k$, and carries
a free $S_k$-action by coordinate permutation. The quotient
$\Conf_k(M)/S_k = UF(k, M)$ is the unordered configuration space.
\end{definition}

\begin{example}[Euclidean case]
\label{ex:conf-Rn}
For $M = \mathbb{R}^n$ the configuration space $\Conf_k(\mathbb{R}^n)$
retracts onto the \emph{reduced} configuration space
\[
 \widetilde{\Conf}_k(\mathbb{R}^n)
 \;=\; \{(x_1, \ldots, x_k) :
        \bar x = 0,\ \sum |x_i - \bar x|^2 = 1,\ x_i \neq x_j\},
\]
a closed manifold of dimension $kn - n - 1$. The cohomology is known:
\[
 H^\bullet(\Conf_k(\mathbb{R}^n); \mathbb{Q})
 \;=\; \mathbb{Q}[g_{ij}]_{1 \leq i < j \leq k}
        /\langle g_{ij}^2 = 0,\ g_{ij} g_{jk} + g_{jk} g_{ki} + g_{ki} g_{ij} = 0 \rangle,
\]
with $|g_{ij}| = n - 1$. The cubic relations are the Arnold
quadratic relations \cite[Arnold 1969]{Arnold1969}. For $n = 3$
the generator is in degree $2$, and the relations produce the
cohomology of the little $3$-disks operad as the free graded
commutative algebra on the $g_{ij}$ modulo Arnold.
\end{example}

\begin{definition}[Fulton--MacPherson compactification]
\label{def:fm-compactification}
Let $M$ be a smooth manifold and $k \geq 1$. The
\emph{Fulton--MacPherson compactification}
$\FMbar(k, M)^{\FM}$ \cite{FultonMacPherson1994} is the closure of
the image of $\Conf_k(M) \hookrightarrow M^k \times
\prod_{|I| \geq 2} \mathrm{Bl}_{\Delta_I}(M^I)$ under the simultaneous
inclusions, where $\Delta_I \subset M^I$ is the thin diagonal indexed
by $I \subseteq \{1, \ldots, k\}$ and $\mathrm{Bl}$ is real blow-up.
\end{definition}

\begin{theorem}[Fulton--MacPherson 1994]
\label{thm:fm-basic}
The space $\FMbar(k, M)^{\FM}$ is a smooth manifold with corners,
compact when $M$ is compact. The interior is
$\Conf_k(M) \subset \FMbar(k, M)^{\FM}$ as an open dense submanifold.
The boundary $\partial \FMbar(k, M)^{\FM}$ is stratified by labelled
rooted trees $T$ with $k$ leaves and no bivalent internal vertices;
the stratum $S_T \subset \partial \FMbar(k, M)^{\FM}$ is canonically
identified with
\[
 S_T \;\simeq\; \Conf_{|v_T|}(M) \times
 \prod_{v \in \mathrm{int}(T)} \Conf_{|v|}(\mathbb{R}^n),
\]
where $v_T$ is the root vertex, $|v_T|$ the number of children of
the root, $\mathrm{int}(T)$ the set of non-root internal vertices,
and $|v|$ the number of children of each internal vertex. The
$\Conf_{|v|}(\mathbb{R}^n)$-factors record the \emph{relative}
scale-and-position data of points colliding into the cluster at
$v$, modulo $\mathbb{R}^n \rtimes \mathbb{R}_{>0}$ (translation and
scaling).
\end{theorem}

The tree-stratification is the operadic engine. Collisions of $k$
points along a rooted tree $T$ correspond to nested sub-configurations:
each internal vertex $v$ of $T$ with children $v_1, \ldots, v_{|v|}$
represents a cluster where the points labelled by the leaves beneath
$v_1, \ldots, v_{|v|}$ collide at the same rate and in a relative
configuration given by a point of $\Conf_{|v|}(\mathbb{R}^n)$.

\begin{example}[Arity-$2$ and arity-$3$ Euclidean cases at $n = 3$]
\label{ex:fm-arity-2-3}
For $M = \mathbb{R}^3$:
\begin{itemize}
 \item $\FMbar(1, \mathbb{R}^3)^{\FM} = \mathbb{R}^3$, trivial.
 \item $\FMbar(2, \mathbb{R}^3)^{\FM} = \mathbb{R}^3 \times
       \mathrm{Bl}_0(\mathbb{R}^3)$, the oriented blow-up of the
       diagonal. The interior is $\Conf_2(\mathbb{R}^3)$; the
       boundary is a copy of $\mathbb{R}^3 \times S^2$ recording the
       unit-direction along which the two points collide. Modulo
       translation this is $S^2 = \Conf_2(\mathbb{R}^3)/\mathbb{R}^3
       \rtimes \mathbb{R}_{>0}$.
 \item $\FMbar(3, \mathbb{R}^3)^{\FM}$: interior
       $\Conf_3(\mathbb{R}^3)$; boundary is the union of three
       codimension-$1$ strata (binary trees with root $\{1,2,3\}$ and
       one nontrivial internal vertex: $\{(12)3\}$, $\{(13)2\}$,
       $\{(23)1\}$, each fibred over a configuration of
       $\binom{3}{2}$ points by a $\Conf_2(\mathbb{R}^3)$ fibre modulo
       scale-and-translate) together with a codimension-$2$ stratum
       (the ternary tree $\{1, 2, 3\}$ colliding simultaneously,
       producing a copy of $\Conf_3(\mathbb{R}^3)/\mathbb{R}^3
       \rtimes \mathbb{R}_{>0}$ of real dimension $6$).
\end{itemize}
The reduced compactification
$\widetilde{\FMbar}(k, \mathbb{R}^n)^{\FM} = \FMbar(k, \mathbb{R}^n)^{\FM}/\mathbb{R}^n \rtimes \mathbb{R}_{>0}$
is the arity-$k$ space of the Fulton--MacPherson operad $\FM_n$.
\end{example}

\begin{heal}[adding chain-level operadic content]
\label{heal:cycle-2}
With Definition~\ref{def:fm-compactification} and
Theorem~\ref{thm:fm-basic} the $E_3$-label acquires \emph{content}:
the arity-$k$ operations are the points of
$\widetilde{\FMbar}(k, \mathbb{R}^3)^{\FM}$; the boundary strata record
operadic composition; the singular chain complex
$C_\bullet(\widetilde{\FMbar}(k, \mathbb{R}^3)^{\FM}; \mathbb{Q})$ is
the chain-level model entering Kontsevich formality.
\end{heal}

%==============================================================================
\section{Cycle 3: ATTACK --- operadic composition is unstated}
%==============================================================================

\begin{attack}[How do FM spaces compose?]
\label{att:cycle-3}
Definition~\ref{def:fm-compactification} produces a sequence of spaces
indexed by arity $k$. An operad requires more: composition maps
\[
 \gamma_k^{\ell_1, \ldots, \ell_k} \colon
 \widetilde{\FMbar}(k, \mathbb{R}^n)^{\FM}
 \times \prod_{i=1}^{k} \widetilde{\FMbar}(\ell_i, \mathbb{R}^n)^{\FM}
 \longrightarrow
 \widetilde{\FMbar}(\ell_1 + \cdots + \ell_k, \mathbb{R}^n)^{\FM},
\]
satisfying associativity and $S_{\ell_i}$-equivariance. Do these exist
for the Fulton--MacPherson compactification? If so, are they associative?
Are they $S_k$-equivariant? Do they restrict to the composition of the
little $n$-disks operad on the interior?
\end{attack}

\begin{survivor}[\ref{att:cycle-3}]
The Fulton--MacPherson compactification supports operadic composition
by \emph{attaching} sub-configurations along the boundary stratification.
The composition map $\gamma$ is a diffeomorphism onto the
codimension-$k - 1$ boundary stratum corresponding to simultaneous
collision of all sub-configurations.
\end{survivor}

\begin{construction}[Operadic composition for $\FM_n$]
\label{constr:fm-composition}
Fix arities $k, \ell_1, \ldots, \ell_k \geq 1$ and write
$L = \ell_1 + \cdots + \ell_k$. Given
$x = (x_1, \ldots, x_k) \in \widetilde{\FMbar}(k, \mathbb{R}^n)^{\FM}$
and sub-configurations
$y^{(i)} = (y^{(i)}_1, \ldots, y^{(i)}_{\ell_i}) \in
 \widetilde{\FMbar}(\ell_i, \mathbb{R}^n)^{\FM}$
for $i = 1, \ldots, k$, the composition
$\gamma(x; y^{(1)}, \ldots, y^{(k)}) \in
 \widetilde{\FMbar}(L, \mathbb{R}^n)^{\FM}$
is the limit configuration
\[
 (x_i + \varepsilon y^{(i)}_j)_{i = 1, \ldots, k,\ j = 1, \ldots, \ell_i}
 \qquad \text{as } \varepsilon \to 0^+,
\]
extended to the compactification by the tree-stratification
of Theorem~\ref{thm:fm-basic}: the limit lives in the boundary
stratum labelled by the rooted $2$-level tree with root arity $k$ and
children arities $\ell_1, \ldots, \ell_k$. This defines a smooth
diffeomorphism
\[
 \gamma_k^{\ell_1, \ldots, \ell_k} \colon
 \widetilde{\FMbar}(k, \mathbb{R}^n)^{\FM}
 \times \prod_{i} \widetilde{\FMbar}(\ell_i, \mathbb{R}^n)^{\FM}
 \xrightarrow{\ \sim\ }
 S_{T_{k, \ell_1, \ldots, \ell_k}}
 \;\subset\;
 \widetilde{\FMbar}(L, \mathbb{R}^n)^{\FM},
\]
where $S_T$ is the codimension-$k - 1$ stratum of
Theorem~\ref{thm:fm-basic}. Associativity for three-fold composition
follows from the fact that both orders assemble the same rooted
$3$-level tree; $S_k$-equivariance is the permutation of labels.
\end{construction}

\begin{theorem}[$\FM_n$ is an operad in smooth manifolds with corners]
\label{thm:fm-operad}
The collection $\FM_n = \{\widetilde{\FMbar}(k, \mathbb{R}^n)^{\FM}\}_{k \geq 1}$
with composition of Construction~\ref{constr:fm-composition} is an
operad in the category $\mathbf{Man}^\partial_\infty$ of smooth
manifolds with corners and smooth maps. It is a cofibrant model for
the little $n$-disks operad $\Ddisk_n$
\cite[Salvatore 2001, Markl--Shnider--Stasheff 2002, Sinha 2004]%
{Salvatore2001, MarklShniderStasheff2002, Sinha2004}: there is a weak
equivalence $\Ddisk_n \xrightarrow{\sim} \FM_n$ of topological operads,
and $\FM_n$ is freely generated by its arity-$1$ and the boundary
attachment data, making it $\Sigma$-cofibrant in the sense of
Berger--Moerdijk.
\end{theorem}

\begin{heal}[FM is an honest operad]
\label{heal:cycle-3}
Construction~\ref{constr:fm-composition} and Theorem~\ref{thm:fm-operad}
exhibit the composition structure explicitly, closing the gap from
Attack~\ref{att:cycle-3}. The little $n$-disks operad $\Ddisk_n$
is the homotopical target; $\FM_n$ is its smooth-manifold model,
suitable for chain-level algebraic manipulation.
\end{heal}

%==============================================================================
\section{Cycle 4: ATTACK --- what is the chain-level $E_n$-operad?}
%==============================================================================

\begin{attack}[Chain model unspecified]
\label{att:cycle-4}
``Chain-level $E_3$-operad'' is itself a free-floating phrase. The
natural candidates --- singular chains $C_\bullet(\FM_n; \mathbb{Q})$,
semi-algebraic chains $C^{\mathrm{sa}}_\bullet(\FM_n)$, polynomial
differential forms $\Omega^\bullet_{\mathrm{PA}}(\FM_n)$ --- are all
quasi-isomorphic but not equal; pinning down which model is operative
is necessary to state Kontsevich formality as a chain-level theorem.
Moreover, the $\FM_n$ operad in topological spaces does not automatically
lift to an operad in chain complexes: operadic composition must be
promoted through a diagonal approximation compatible with the boundary
stratification.
\end{attack}

\begin{survivor}[\ref{att:cycle-4}]
The correct chain-level model is the
\emph{semi-algebraic chain complex}
$C^{\mathrm{sa}}_\bullet(\FM_n; \mathbb{Q})$ of Kontsevich--Soibelman
\cite[KS 2000, \S 3]{KontsevichSoibelman2000-deformations},
or equivalently the polynomial de Rham model
$\Omega^\bullet_{\mathrm{PA}}(\FM_n)$ of Hardt--Lambrechts--Turchin--Volic
\cite[HLTV 2011]{HLTV2011}. Both models are strictly operadic in chain
complexes over $\mathbb{Q}$ and are strictly multiplicative under
operadic composition, avoiding the Eilenberg--Zilber correction needed
for singular chains.
\end{survivor}

\begin{definition}[Semi-algebraic chain operad]
\label{def:sa-chain-op}
For a compact semi-algebraic manifold with corners $X$ let
$C^{\mathrm{sa}}_\bullet(X; \mathbb{Q})$ denote the complex of
semi-algebraic chains with rational coefficients (Kontsevich--Soibelman
loc.~cit.): a chain is a formal $\mathbb{Q}$-linear combination of
compact oriented semi-algebraic subsets modulo semi-algebraic
equivalence, with boundary given by the topological boundary with
induced orientation. Cartesian products of semi-algebraic chains are
strictly multiplicative: $C^{\mathrm{sa}}_\bullet(X \times Y)
\supset C^{\mathrm{sa}}_\bullet(X) \otimes C^{\mathrm{sa}}_\bullet(Y)$
is a strict inclusion of chain complexes. Applied to
$\FM_n$ term-by-term:
\[
 C_n := C^{\mathrm{sa}}_\bullet(\FM_n; \mathbb{Q})
 \;=\; \{C^{\mathrm{sa}}_\bullet(\widetilde{\FMbar}(k, \mathbb{R}^n)^{\FM}; \mathbb{Q})\}_{k \geq 1},
\]
with operadic composition induced from Construction~\ref{constr:fm-composition}
by the strict Cartesian multiplicativity of semi-algebraic chains:
$\gamma_*(\sigma \otimes \tau_1 \otimes \cdots \otimes \tau_k) =
 \gamma_!(\sigma \times \tau_1 \times \cdots \times \tau_k)$
where $\gamma_!$ is the push-forward along the smooth diffeomorphism onto
the boundary stratum. This is a dg-operad in $\mathrm{Ch}_\mathbb{Q}$.
\end{definition}

\begin{proposition}[Weak equivalence between chain models]
\label{prop:chain-models-equivalent}
The three candidate chain-level models for $E_n$,
\[
 C^{\mathrm{sa}}_\bullet(\FM_n; \mathbb{Q}),
 \quad
 \Omega^\bullet_{\mathrm{PA}}(\FM_n; \mathbb{Q})^\vee,
 \quad
 C_\bullet(\Ddisk_n; \mathbb{Q}),
\]
are pairwise quasi-isomorphic as dg-operads over $\mathbb{Q}$. The
first two are strictly operadic; the third is operadic only up to a
coherent Eilenberg--Zilber homotopy
\cite[HLTV 2011 \S 5; Fresse 2017]{HLTV2011, Fresse2017}.
\end{proposition}

\begin{theorem}[Cohomology operad of $\FM_n$, after F.~Cohen]
\label{thm:fm-cohomology}
For $n \geq 1$ the cohomology operad of $\FM_n$ is
\[
 H^\bullet(\FM_n; \mathbb{Q}) \;=\;
 \begin{cases}
  \mathrm{Ass}, & n = 1, \\
  \Gerst_n, & n \geq 2,
 \end{cases}
\]
where $\Gerst_n$ is the operad of $(n-1)$-shifted Poisson algebras:
free graded commutative product $\cdot$ in degree $0$ plus a Lie
bracket $\{-, -\}$ in degree $-(n-1)$, satisfying Leibniz. For $n = 3$
the bracket has degree $-2$. The generators of $H^\bullet(\FM_n(k);
\mathbb{Q})$ are the pairwise Arnold classes $g_{ij}$ in degree $n-1$
(Example~\ref{ex:conf-Rn}); the quadratic relations recover
the Gerstenhaber axioms.
\end{theorem}

\begin{heal}[chain model fixed]
\label{heal:cycle-4}
Definition~\ref{def:sa-chain-op}, Proposition~\ref{prop:chain-models-equivalent},
and Theorem~\ref{thm:fm-cohomology} pin down the chain-level $E_n$-operad
as $C_n = C^{\mathrm{sa}}_\bullet(\FM_n; \mathbb{Q})$, strictly operadic,
with cohomology $\Gerst_n$ for $n \geq 2$. Kontsevich formality will
be the statement that $C_n$ and $\Gerst_n$ are quasi-isomorphic as
dg-operads, up to a torsor action.
\end{heal}

%==============================================================================
\section{Cycle 5: ATTACK --- where does $\GRT_1$ come from?}
%==============================================================================

\begin{attack}[Torsor structure unmotivated]
\label{att:cycle-5}
Cached statement from the cache-layer: ``Stage-$1$ canonical up to
$\GRT_1(\mathbb{Q})$-torsor''. This assertion makes sense only if there
is a well-defined space of $\FM_n$-formality morphisms
$C_n \to \Gerst_n$, and this space is a $\GRT_1(\mathbb{Q})$-torsor.
The chapter states the torsor but does not exhibit any of:
\begin{itemize}
 \item the formality morphism as a $\GRT_1$-torsor element;
 \item the identification of $\GRT_1$ with a graph-complex
       cohomology group;
 \item the mechanism by which associator choices and graph-integral
       weights transfer.
\end{itemize}
\end{attack}

\begin{survivor}[\ref{att:cycle-5}]
Tamarkin \cite{Tamarkin2003}, Kontsevich \cite{Kontsevich2003},
Kontsevich--Soibelman \cite{KontsevichSoibelman2000-deformations},
and Willwacher \cite{Willwacher2015} together exhibit the complete
structure: an explicit $\GRT_1$-action on the space of
$E_n$-formality morphisms, realised via Kontsevich graph-complex
cohomology and the integral-weight construction on
$\FMbar(k, \mathbb{R}^n)^{\FM}$.
\end{survivor}

\begin{definition}[Kontsevich graph complex $\GC_n$]
\label{def:graph-complex}
Fix $n \geq 2$. The \emph{Kontsevich graph complex} $\GC_n$ is the dg
vector space (over $\mathbb{Q}$) spanned by isomorphism classes of
connected graphs $\Gamma$ with:
\begin{itemize}
 \item no tadpoles or parallel edges;
 \item all internal vertices of valence $\geq 3$;
 \item each vertex equipped with a total cyclic ordering of the incident
       half-edges if $n$ is even, or no ordering if $n$ is odd.
\end{itemize}
The degree of $\Gamma$ is
\[
 |\Gamma| \;=\; (n - 1) E(\Gamma) - n V(\Gamma) + n,
\]
where $V, E$ are numbers of vertices and edges. The differential
$d_{\GC_n}$ contracts a single edge in all possible ways, with signs
determined by the cyclic ordering. The cohomology $H^\bullet(\GC_n)$
is graded by degree, with a piece in degree zero.
\end{definition}

\begin{theorem}[Willwacher 2015]
\label{thm:willwacher}
There is an isomorphism of graded Lie algebras
\[
 H^0(\GC_2) \;\simeq\; \mathfrak{grt}_1,
\]
the Grothendieck--Teichm\"uller Lie algebra
\cite[Willwacher 2015, \emph{Invent.\ math.}\ 200, Thm.\ 1.1]{Willwacher2015}.
For $n \geq 3$ there is a (Bott-periodicity) identification
\[
 H^0(\GC_n) \;\simeq\; \mathfrak{grt}_1 \quad (n \geq 2),
\]
realising $\mathfrak{grt}_1$ as a universal piece of graph cohomology
independent of dimension. Exponentiating produces
$\GRT_1(\mathbb{Q}) = \exp(\mathfrak{grt}_1)$ acting on all $E_n$-operad
models uniformly.
\end{theorem}

\begin{construction}[Kontsevich graph-integral operation]
\label{constr:graph-integral}
For a graph $\Gamma$ with $k$ \emph{external} (leaf) vertices and $m$
\emph{internal} vertices, the \emph{Kontsevich integral weight}
$W_\Gamma$ is the integral over $\widetilde{\FMbar}(k + m, \mathbb{R}^n)^{\FM}$
of a product of $2$-forms, one for each edge, pulled back from
$\widetilde{\FMbar}(2, \mathbb{R}^n)^{\FM} \simeq S^{n-1}$ via the
forgetful maps that remember only the two endpoints of each edge
\cite[Kontsevich 2003, \S 6.4]{Kontsevich2003}. The integrand is a
top-degree closed form on the corner manifold; convergence follows
from Fulton--MacPherson smoothness and Stokes' theorem.
\end{construction}

\begin{theorem}[Kontsevich--Tamarkin $E_n$-formality, torsor structure]
\label{thm:kt-formality}
For $n \geq 2$ there exists a quasi-isomorphism of dg-operads
$\mathcal{F} \colon \Gerst_n \xrightarrow{\sim} C_n$, called an
$E_n$-formality morphism. The space of such morphisms, up to
homotopy, is a (non-empty) torsor under the group
$\GRT_1(\mathbb{Q})$, acting through the graph-complex realisation
of Theorem~\ref{thm:willwacher}: the image of a formality morphism
under $\xi \in \GRT_1(\mathbb{Q})$ is the formality morphism whose
graph-integral weights $W_\Gamma$ are twisted by $\xi$ on each graph
$\Gamma \in \GC_n$.
\end{theorem}

\begin{heal}[$\GRT_1$ acts on formality morphisms]
\label{heal:cycle-5}
Theorem~\ref{thm:kt-formality} closes Attack~\ref{att:cycle-5}: the
$\GRT_1$-torsor structure on Stage-$1$ formality is the torsor of
graph-integral weights twisted by $H^0(\GC_n) = \mathfrak{grt}_1$.
In the Vol~III correspondence $\Phi_3 = \SpCh_{\Sigma_2, C} \circ
\PhiFA_3$, different choices in the $\GRT_1$-torsor produce
Stage-$1$ outputs that are $E_3$-quasi-isomorphic but differ by
a super-KZ-type associator on the chain level; the difference is
absorbed into the Stage-$2$ specialisation datum $(\Sigma_2, C)$ and
does not affect the output chiral algebra up to isomorphism.
\end{heal}

%==============================================================================
\section{Cycle 6: ATTACK --- how does $\PhiFA_3$ \emph{use} $E_3$?}
%==============================================================================

\begin{attack}[Stage-$1$ assembly unstated at chain level]
\label{att:cycle-6}
Granted the explicit $E_3$ operad $C_3 = C^{\mathrm{sa}}_\bullet(\FM_3; \mathbb{Q})$
and Kontsevich formality $\mathcal{F} \colon \Gerst_3 \to C_3$, it
remains to say how $\PhiFA_3$ uses these operations. The chapter
asserts $\PhiFA_3$ produces an ``$E_3$-holomorphic factorisation
algebra'' on $X$ from the degree-$(-2)$ Gerstenhaber bracket on
$\HH^\bullet(\cC)$; the explicit chain-level assignment of
local-observable data to discs remains opaque.
\end{attack}

\begin{survivor}[\ref{att:cycle-6}]
The Stage-$1$ output is a Costello--Gwilliam factorisation algebra
on $X$ in the sense of \cite{CostelloGwilliam2017, CostelloGwilliam2021},
with disc observables assembled from $E_3$ compositions indexed by
$\FMbar(k, \mathbb{R}^3)^{\FM}$ strata. The holomorphic refinement
is local-analytic: open polydiscs in $X$ serve as the local
configuration-space inputs.
\end{survivor}

\begin{definition}[Holomorphic little $3$-discs subcategory]
\label{def:hol-disks}
Fix a smooth complex $\CY_3$ variety $X$. A \emph{holomorphic
$3$-disc} in $X$ is an open embedding
$\varphi \colon B^6_\epsilon \hookrightarrow X$ where $B^6_\epsilon
\subset \mathbb{C}^3$ is the open polydisc of polyradius $\epsilon$,
and $\varphi$ is holomorphic. Two discs $\varphi_i, \varphi_j$ are
\emph{disjoint} if their images have empty intersection. The
category $\mathbf{Disc}_3^{\mathrm{hol}}(X)$ has:
\begin{itemize}
 \item objects: finite collections of pairwise disjoint holomorphic
       $3$-discs in $X$;
 \item morphisms: inclusions that embed one collection into another,
       with the images of the smaller collection contained in the
       images of the larger.
\end{itemize}
\end{definition}

\begin{theorem}[$E_3$ on holomorphic discs via configuration forgetting]
\label{thm:fm-to-disks}
There is a natural functor of topological operads
\[
 \FM_3 \longrightarrow \mathbf{Disc}_3^{\mathrm{hol}}(X)
\]
assigning each configuration
$(x_1, \ldots, x_k) \in \widetilde{\FMbar}(k, \mathbb{R}^3)^{\FM}$
(with chosen polyradius $\epsilon$ small enough for pairwise disjointness)
the collection of holomorphic polydiscs centred at the $x_i$, transported
by a holomorphic chart of $X$. Operadic composition commutes with
disc insertion \cite[Costello--Gwilliam 2021, Ch.\ 3]{CostelloGwilliam2021}.
\end{theorem}

\begin{construction}[Stage-$1$ $\PhiFA_3(\cC)$ as $\FM_3$-colimit]
\label{constr:phifa3-via-fm}
Given a smooth proper cyclic $\Ainf$ $\CY_3$ category $\cC$, the
Stage-$1$ output $\PhiFA_3(\cC)$ is a factorisation algebra on $X$
whose value on a single holomorphic polydisc $B \subset X$ is
$\HH_\bullet(\cC)$ with its $E_3$-algebra structure induced by
Kontsevich--Tamarkin formality; the chain-level $E_3$ action on
$\HH_\bullet(\cC)$ is through the chain operad $C_3$, and the
structure maps on disjoint unions of discs are recovered from
operadic composition on $\FM_3$ via Theorem~\ref{thm:fm-to-disks}.
Concretely, for $k$ disjoint discs $B_1, \ldots, B_k \subset X$:
\[
 \PhiFA_3(\cC)(B_1 \sqcup \cdots \sqcup B_k)
 \;\simeq\; C^{\mathrm{sa}}_\bullet\!\left(\widetilde{\FMbar}(k, \mathbb{R}^3)^{\FM}\right)
  \otimes_{C_3} \HH_\bullet(\cC)^{\otimes k},
\]
the homotopy coend over $\FM_3$ \cite[Costello--Gwilliam 2021, Prop.\ 3.4.3]{CostelloGwilliam2021}.
The Dolbeault operator $\bar\partial$ on $X$ enters through the
local-analytic refinement: the factorisation algebra is a sheaf of
chain complexes with $\bar\partial$-acyclic extensions, producing
the holomorphic refinement of a topological $E_3$-factorisation
algebra \cite[Costello--Li 2016, 2020]{CostelloLi2016, CostelloLi2020}.
\end{construction}

\begin{heal}[chain-level Stage-$1$ made explicit]
\label{heal:cycle-6}
Construction~\ref{constr:phifa3-via-fm} exhibits $\PhiFA_3(\cC)$
explicitly as a $\FM_3$-colimit of $\HH_\bullet(\cC)^{\otimes k}$
assemblies. The operation of $\PhiFA_3$ is now a transparent
assignment: $\cC \mapsto$ cyclic Hochschild chains with
Kontsevich-formality $E_3$-structure, reassembled on discs of $X$
by Fulton--MacPherson operadic composition.
\end{heal}

%==============================================================================
\section{Cycle 7: RE-ATTACK --- $\bar\partial$-refinement}
%==============================================================================

\begin{attack}[Holomorphic refinement specifies what data?]
\label{att:cycle-7}
The chapter distinguishes three levels in the Stage-$1$ output:
\begin{enumerate}
 \item\label{it:att7-a} $E_3$-algebra on a single point (Hochschild cochain complex);
 \item\label{it:att7-b} $E_3$-factorisation algebra on the topological space $|X|$ (disc combinatorics, no complex structure);
 \item\label{it:att7-c} $E_3$-holomorphic factorisation algebra on $X$ (Dolbeault operator, $\Omega_X$, BV rigidity).
\end{enumerate}
The Fulton--MacPherson $\FM_3$ construction above produces~\ref{it:att7-b}.
What additional chain-level data upgrades this to~\ref{it:att7-c}?
The answer must appear at the level of the configuration-space model
itself to be chain-explicit, not as a black-box referral to BV
theory.
\end{attack}

\begin{survivor}[\ref{att:cycle-7}]
The holomorphic refinement adds \emph{two} pieces of chain-level data
to the $\FM_3$-model:
\begin{itemize}
 \item[(i)] A Dolbeault complex $(\Omega^{0, \bullet}(X), \bar\partial)$
       acting on disc observables by sheaf-level derivation.
 \item[(ii)] A holomorphic-volume trivialisation $\Omega_X \in
       H^0(X, K_X)$ providing a preferred scalar against which
       graph-integrals $W_\Gamma$ of Construction~\ref{constr:graph-integral}
       become \emph{actual numbers} (via integration over
       $\widetilde{\FMbar}$).
\end{itemize}
Both pieces enter the Costello--Li quantisation of holomorphic
Chern--Simons on $X$ \cite{CostelloLi2016, CostelloLi2020}.
\end{survivor}

\begin{definition}[Dolbeault-refined $\FM_3$-factorisation algebra]
\label{def:dolbeault-fm}
A \emph{Dolbeault-refined $E_3$-factorisation algebra} on $X$ is a
factorisation algebra $\cA$ in the category of $\bar\partial$-complexes
(Dolbeault dg-sheaves) on $X$, such that:
\begin{itemize}
 \item on an open polydisc $B \subset X$, the value $\cA(B)$ is a
       finite-dimensional cohomology summand of a free-module resolution
       of an $E_3$-algebra (with $\bar\partial = 0$ on the
       representative);
 \item for disjoint discs $B_1, \ldots, B_k$ the factorisation
       product is the $\FM_3$-assembly of
       Construction~\ref{constr:phifa3-via-fm} \emph{twisted} by the
       holomorphic-volume graph-integrals $W_\Gamma$ evaluated on
       $\widetilde{\FMbar}(k, \mathbb{C}^3)^{\FM}$.
\end{itemize}
\end{definition}

\begin{theorem}[Costello--Li 2016 / 2020 refinement]
\label{thm:cl-refinement}
The Stage-$1$ output $\PhiFA_3(\cC)$ of Construction~\ref{constr:phifa3-via-fm}
is naturally a Dolbeault-refined $E_3$-factorisation algebra in the
sense of Definition~\ref{def:dolbeault-fm}, with the trivialisation
$\Omega_X$ entering as the BV-rigidifier of the $\bar\partial$-quasi-free
resolution. The Kontsevich graph-integral weights $W_\Gamma$ of
Construction~\ref{constr:graph-integral}, evaluated on the holomorphic
variant $\widetilde{\FMbar}(k, \mathbb{C}^3)^{\FM}$, compute the
Costello--Li propagators of hCS on $X$
\cite[Costello--Li 2020, \S 3]{CostelloLi2020}.
\end{theorem}

\begin{remark}[Relation to 6d hCS]
\label{rem:cy3-6d-hcs}
The Stage-$1$ $\PhiFA_3(\cC)$ for $\cC = \mathrm{Perf}(X)$ coincides,
up to formality choice, with the observables of 6D holomorphic
Chern--Simons on $X$ with gauge Lie algebra $\mathfrak{g}_X =
\HH^\bullet(\cC)$. This is the Costello--Francis--Gwilliam factorisation
algebra of 5D hCS on $X \times \mathbb{R}_t$ after $\mathbb{R}_t$-compactification,
whose holomorphic sector on $X$ carries the Fulton--MacPherson
$\FM_3$-model. The $\GRT_1$-torsor of Theorem~\ref{thm:kt-formality}
acts on hCS propagators through the Costello--Li scheme-dependence,
preserving observables on nose but rotating the BV-exact
counter-terms \cite{CostelloLi2020, CFG2026}.
\end{remark}

\begin{heal}[holomorphic refinement via graph-integrals on $\widetilde{\FMbar}(k, \mathbb{C}^3)^{\FM}$]
\label{heal:cycle-7}
Theorem~\ref{thm:cl-refinement} and Remark~\ref{rem:cy3-6d-hcs} close
Attack~\ref{att:cycle-7}: the Dolbeault refinement is not a black box
but a specific chain-level sheafification on
$\widetilde{\FMbar}(k, \mathbb{C}^3)^{\FM}$ through Kontsevich
graph-integrals with the holomorphic-volume trivialisation entering
as the integration data.
\end{heal}

%==============================================================================
\section{Cycle 8: RE-ATTACK --- compatibility with Stage-$2$}
%==============================================================================

\begin{attack}[Stage-$2$ forgets $E_3$?]
\label{att:cycle-8}
The native operadic level of Stage-$1$ is $E_3$ (holomorphic), but the
Stage-$2$ specialisation-then-restriction to a curve $C$ yields an
$\Eone$-chiral algebra, by Proposition~\texttt{native-en-level} in the
target chapter. Does the chain-level $\FM_3$-model produced by
$\PhiFA_3$ degenerate coherently under Stage-$2$, or is there a lossy
projection that destroys the $\FM_3$-data?
\end{attack}

\begin{survivor}[\ref{att:cycle-8}]
Factorisation homology along the $(d-1)$-cycle $\Sigma_{d-1}$ is
an $(\infty, 1)$-functor
$\EdHolFA(X) \to \mathbf{Hol}^1\text{-}\mathrm{FA}(X \setminus \Sigma_{d-1})$
(Lurie \emph{HA}~$5.5.3.6$) that reduces the operadic level by exactly
$d - 1$: the assembly on a topological $(d-1)$-cycle forgets $E_{d-1}$
directions and leaves one residual dimension, producing
$E_{d - (d-1)} = \Eone$ on the complement $X \setminus \Sigma_{d-1}$
\cite[Francis--Gaitsgory 2012, Lemma 3.3.4]{FrancisGaitsgory2012}.
At $d = 3$ this is factorisation homology over $\Sigma_2$ forgetting
two Fulton--MacPherson directions and keeping the one direction
residual on the curve $C$.
\end{survivor}

\begin{construction}[$\FM_3 \to \FM_1$ factorisation-homology projection]
\label{constr:fm3-to-fm1}
Fix $\Sigma_2 \subset X$ a smooth real $2$-cycle and $C \subset
X \setminus \Sigma_2$ a holomorphic reference curve. The factorisation
homology integral
\[
 \int_{\Sigma_2} \colon \FM_3\text{-alg} \longrightarrow \FM_1\text{-alg on } C
\]
is implemented at the chain level by the operadic push-forward along
the fibration
\[
 \widetilde{\FMbar}(k, \mathbb{R}^3)^{\FM}
  \;\longrightarrow\;
  \widetilde{\FMbar}(k, \mathbb{R}^1)^{\FM}
\]
obtained by projecting each configuration point $x_i = (x_i^1, x_i^2, x_i^3)$
onto its first coordinate $x_i^1$ and compactifying the pairwise
collisions in $\mathbb{R}^1$. The fibre of this projection over a
configuration $(x_1^1, \ldots, x_k^1) \in \widetilde{\FMbar}(k,
\mathbb{R}^1)^{\FM}$ is a compactification of
$\prod_i \{(x_i^2, x_i^3)\}$ of real dimension $2k$, which is
topologically equivalent to a product of $2$-cells $\{(y_i^2, y_i^3) :
|y_i^2|^2 + |y_i^3|^2 \leq \epsilon^2\}$, after blowing down the
compactification in the normal direction. Factorisation homology
corresponds to the push-forward $\pi_!$ along this fibration,
implemented at the chain level by integration of the $\FM_3$-structure
against the top form on the $(2k)$-dimensional fibres.
\end{construction}

\begin{proposition}[Residual $\Eone$-chiral structure after Stage-$2$]
\label{prop:e1-residual}
For $\cC$ a smooth proper cyclic $\Ainf$ $\CY_3$ category and
$(\Sigma_2, C)$ an admissible specialisation datum in $X$, the
composite
\[
 \Phi_3^{(\Sigma_2, C)}(\cC) \;=\;
  \SpCh_{\Sigma_2, C}\!\left(\PhiFA_3(\cC)\right)
\]
is an $\Eone$-chiral algebra on $C$ in the sense of Beilinson--Drinfeld
chiral algebras, with operadic operations indexed by the residual
$\FM_1$-model of Construction~\ref{constr:fm3-to-fm1} composed with the
holomorphic pullback to $C$. The $\pi_0(\Conf_2(\mathbb{R}^3))$-symmetry
of the $E_3$-level is lost in this residual: the $E_3 \to \Etwo$
restriction via Dunn additivity produces a \emph{symmetric} braiding
(since $\pi_1(\Conf_2(\mathbb{R}^3)) = 0$ by
Fadell--Neuwirth~\cite{FadellNeuwirth1962}), and the
$\Etwo \to \Eone$ further restriction by $\Sigma_2$-integration forgets
the symmetric braiding entirely. The non-symmetric $R$-matrix data,
where it exists, is recovered from the Drinfeld centre
$\cZ(\mathrm{Rep}^{\Eone}(A_\cC))$, not from the chain-level
$\FM_3$-structure \cite[target chapter \texttt{cy\_to\_chiral.tex}
Theorem~\texttt{c3-drinfeld-center}]{VolIII}.
\end{proposition}

\begin{heal}[chain-level Stage-$2$ projection is explicit]
\label{heal:cycle-8}
Construction~\ref{constr:fm3-to-fm1} and Proposition~\ref{prop:e1-residual}
close Attack~\ref{att:cycle-8}: the Stage-$2$ transition from the
chain-level $\FM_3$-model to the $\FM_1$-model on $C$ is an explicit
chain-level push-forward along the projection fibration, forgetting
exactly the $\Sigma_{d-1} = \Sigma_2$ directions of the configuration
space. The operadic-level reduction $E_3 \to \Eone$ is a
chain-explicit degeneration, not a lossy projection at the level of
$\FM_3$ data.
\end{heal}

%==============================================================================
\section{Cycle 9: RE-ATTACK --- coherence with Dunn--Lurie additivity}
%==============================================================================

\begin{attack}[Dunn--Lurie at the chain level]
\label{att:cycle-9}
Dunn additivity \cite[Dunn 1988; Lurie \emph{HA} 5.1.2.2]{Dunn1988, LurieHA}
asserts $E_m \otimes E_n \simeq E_{m+n}$ as operads in a symmetric
monoidal $(\infty, 1)$-category. In the Vol~III chapter this is used
to derive obstruction-theoretic identities like
$\HH^{-2}_{E_1}(\HH_\bullet(\cC)) = 0$ (Theorem~\texttt{derived-framing-obstruction})
and the symmetric-vs-non-symmetric distinction of the $\Etwo$-braiding.
At the chain level of $\FM_3$, the corresponding statement is
$\FM_3 \simeq \FM_2 \otimes \FM_1$; this is a statement about manifolds
with corners that requires explicit substantiation. Is the isomorphism
a diffeomorphism of operads in smooth manifolds with corners, or only
a quasi-isomorphism of chain operads?
\end{attack}

\begin{survivor}[\ref{att:cycle-9}]
At the level of topological spaces, the Dunn--Lurie additivity
$\FM_m \otimes \FM_n \simeq \FM_{m+n}$ is a weak equivalence of
operads in topological spaces (Boardman--Vogt 1973), not a
diffeomorphism of manifolds with corners. The smooth-manifold
structure on $\FMbar(k, \mathbb{R}^m \times \mathbb{R}^n)^{\FM}$
is inherited from the global configuration space
$\Conf_k(\mathbb{R}^m \times \mathbb{R}^n)$, which is
diffeomorphic-with-boundary to neither a Cartesian product of
$\FMbar(k, \mathbb{R}^m)^{\FM}$ nor a tensor-product mixture.
The equivalence of $\FM_3$ with $\FM_2 \otimes \FM_1$ is therefore a
homotopical statement at the level of chain complexes: the
chain-level dg-operad $C^{\mathrm{sa}}_\bullet(\FM_3; \mathbb{Q})$ is
quasi-isomorphic to $C^{\mathrm{sa}}_\bullet(\FM_2; \mathbb{Q}) \otimes
C^{\mathrm{sa}}_\bullet(\FM_1; \mathbb{Q})$ as dg-operads in
$\mathrm{Ch}_\mathbb{Q}$.
\end{survivor}

\begin{theorem}[Chain-level Dunn additivity for $\FM_n$]
\label{thm:fm-dunn}
For $m, n \geq 1$ there is a quasi-isomorphism of dg-operads in
$\mathrm{Ch}_\mathbb{Q}$
\[
 C^{\mathrm{sa}}_\bullet(\FM_{m+n}; \mathbb{Q})
 \;\xrightarrow{\ \sim\ }\;
 C^{\mathrm{sa}}_\bullet(\FM_m; \mathbb{Q}) \otimes_{H}
 C^{\mathrm{sa}}_\bullet(\FM_n; \mathbb{Q}),
\]
where $\otimes_H$ denotes the Hadamard (Boardman--Vogt) tensor
product of operads \cite[Boardman--Vogt 1973; Fresse 2017]{BoardmanVogt1973, Fresse2017}.
The quasi-isomorphism is given at the arity-$k$ level by the projection
\[
 \widetilde{\FMbar}(k, \mathbb{R}^{m+n})^{\FM}
 \;\longrightarrow\;
 \widetilde{\FMbar}(k, \mathbb{R}^m)^{\FM} \times
 \widetilde{\FMbar}(k, \mathbb{R}^n)^{\FM},
 \quad
 (x^1_i, \ldots, x^{m+n}_i) \mapsto
 \left((x^1_i, \ldots, x^m_i),\ (x^{m+1}_i, \ldots, x^{m+n}_i)\right),
\]
with compactification extended boundary-component-wise.
\end{theorem}

\begin{corollary}[Compatibility with target-chapter $E_3$ obstruction theorem]
\label{cor:compat-e3-obstr}
Theorem~\ref{thm:fm-dunn} at $(m, n) = (2, 1)$ recovers the Dunn
additivity $\FM_3 \simeq \FM_2 \otimes \FM_1$ used in the target
chapter's derivation of
$\HH^{-2}_{E_1}(\HH_\bullet(\cC)) = 0$: the $E_3$-to-$\Etwo$
forgetful functor restricts to an $\Eone$-algebra on
$\HH_\bullet(\cC)$ in the category of $\Etwo$-algebras, and the
obstruction for promoting $\Etwo$ to $E_3$ is $\HH^{-2}_{E_1}$ of
$\Etwo$-algebras, which is the chain-level translation of the
$\FM_2 \otimes \FM_1$-decomposition of $\FM_3$ at the tangent-complex
level.
\end{corollary}

\begin{heal}[Dunn additivity is chain-homotopical]
\label{heal:cycle-9}
Theorem~\ref{thm:fm-dunn} and Corollary~\ref{cor:compat-e3-obstr}
close Attack~\ref{att:cycle-9}: Dunn additivity for $\FM_3$ is a
quasi-isomorphism of chain operads, and the target-chapter
obstruction-theoretic computations are chain-level shadows of this
quasi-isomorphism.
\end{heal}

%==============================================================================
\section{Cycle 10: RE-ATTACK --- Boardman--Vogt $W$-construction}
%==============================================================================

\begin{attack}[Cofibrancy and $W$-construction]
\label{att:cycle-10}
Theorem~\ref{thm:fm-operad} asserts $\FM_n$ is a cofibrant model for
$\Ddisk_n$, used implicitly to justify that $\FM_n$-algebras and
$\Ddisk_n$-algebras coincide homotopically. The cofibrancy is
through the Boardman--Vogt $W$-construction
\cite{BoardmanVogt1973, BergerMoerdijk2006}, which produces a
canonical cofibrant resolution of any topological operad. The
chapter does not exhibit the comparison; the intervention of the
$W$-construction deserves a chain-level statement.
\end{attack}

\begin{survivor}[\ref{att:cycle-10}]
The Boardman--Vogt $W$-construction $W P$ of a topological operad
$P$ replaces $P$ by its ``trees with length-decorated edges''. For
$P = \Ddisk_n$ the little $n$-discs operad, $W \Ddisk_n$ is
weakly equivalent to $\FM_n$ via a tree-level map that identifies
each length-decorated tree with the boundary stratum it parametrises
in $\FMbar(k, \mathbb{R}^n)^{\FM}$. This makes $\FM_n$ a \emph{small}
cofibrant resolution of $\Ddisk_n$, suitable for chain-level algebra.
\end{survivor}

\begin{definition}[$W$-construction]
\label{def:w-construction}
For a topological operad $P$ with arity-$k$ space $P(k)$, the
\emph{Boardman--Vogt $W$-construction} $W P$ has arity-$k$ space
the moduli of isomorphism classes of pairs $(T, \vec\ell)$ where:
\begin{itemize}
 \item $T$ is a rooted tree with $k$ leaves;
 \item $\vec\ell$ is an assignment of a length $\ell_e \in [0, \infty]$
       to each internal edge $e$;
 \item each internal vertex $v$ of $T$ is decorated by a point of
       $P(|v|)$,
\end{itemize}
modulo the relations: $(T, \vec\ell)$ with $\ell_e = 0$ is identified
with the tree obtained by contracting $e$ (vertex composition via the
operadic structure on $P$); $(T, \vec\ell)$ with $\ell_e = \infty$
is identified with the tree obtained by inserting a unit at the
endpoint of $e$. The operadic composition on $W P$ is tree-grafting
with the new edge receiving length $\infty$ (inclusion of the operad
unit).
\end{definition}

\begin{theorem}[Berger--Moerdijk comparison]
\label{thm:bm-comparison}
\cite[Berger--Moerdijk 2006, Thm 5.1]{BergerMoerdijk2006}
The canonical map $W P \to P$ obtained by collapsing all edge lengths
to zero is a weak equivalence of topological operads. $W P$ is
$\Sigma$-cofibrant. For $P = \Ddisk_n$ the little $n$-discs, there is
a weak equivalence of operads
\[
 W \Ddisk_n \xrightarrow{\ \sim\ } \FM_n,
\]
identifying a length-decorated rooted tree with $k$ leaves
$(T, \vec\ell, \vec p)$ with the boundary stratum of
$\FMbar(k, \mathbb{R}^n)^{\FM}$ labelled by the same tree, where the
length $\ell_e$ records the relative scale of the children cluster
at vertex $v$ and the decoration $\vec p \in \prod_v P(|v|)$ records
their relative angular position.
\end{theorem}

\begin{remark}[$\FM_n$ is minimal; $W \Ddisk_n$ is canonical]
\label{rem:fm-minimal}
$\FM_n$ is a cofibrant model, but Theorem~\ref{thm:bm-comparison}
does not make it \emph{the} cofibrant replacement; the $W$-construction
$W \Ddisk_n$ is canonical as a left-adjoint-type resolution. For
chain-level computations (e.g., formality, graph-integrals), $\FM_n$
is the preferred model because it is finite-dimensional
$(2k + n - 1)$-dimensional at arity $k$ and has smooth-manifold
corners \cite[Sinha 2006]{Sinha2006}.
\end{remark}

\begin{heal}[cofibrancy made manifest]
\label{heal:cycle-10}
Theorem~\ref{thm:bm-comparison} and Remark~\ref{rem:fm-minimal}
close Attack~\ref{att:cycle-10}: $\FM_n$ is a cofibrant replacement
for $\Ddisk_n$ via $W \Ddisk_n$; the chain-level algebra operates
on $\FM_n$ (not $\Ddisk_n$) precisely to exploit the smooth-manifold
cofibrancy.
\end{heal}

%==============================================================================
\section{Summary: surviving core after 10 attack--heal cycles}
%==============================================================================

\begin{survivor}[Concrete $E_3$ on the output side of $\PhiFA_3$]
\label{sur:core}
After ten attack--heal cycles, the following pieces of $E_3$
structure on the output side of $\PhiFA_3$ have stable chain-level
content (no outstanding attacks):
\begin{enumerate}[label=(S\arabic*), leftmargin=*]
 \item\label{S1} The arity-$k$ space of operations is the
      Fulton--MacPherson compactified configuration space
      $\widetilde{\FMbar}(k, \mathbb{R}^3)^{\FM}$, a smooth manifold
      with corners of real dimension $3k - 4$ (Theorem~\ref{thm:fm-basic};
      Example~\ref{ex:fm-arity-2-3}).
 \item\label{S2} Operadic composition is boundary attachment via
       Construction~\ref{constr:fm-composition}, associative and
       $S_k$-equivariant (Theorem~\ref{thm:fm-operad}).
 \item\label{S3} The chain-level model is
       $C_3 = C^{\mathrm{sa}}_\bullet(\FM_3; \mathbb{Q})$ (Definition~\ref{def:sa-chain-op}),
       quasi-isomorphic to $\Ddisk_3$ chains (Proposition~\ref{prop:chain-models-equivalent}).
 \item\label{S4} The cohomology operad is $\Gerst_3$
       (Theorem~\ref{thm:fm-cohomology}): degree-$0$ commutative product,
       degree-$(-2)$ Lie bracket.
 \item\label{S5} Kontsevich $E_3$-formality $\mathcal{F} \colon
       \Gerst_3 \xrightarrow{\sim} C_3$ exists and is unique up to a
       $\GRT_1(\mathbb{Q})$-torsor (Theorem~\ref{thm:kt-formality}).
 \item\label{S6} The torsor is realised by Willwacher's
       $H^0(\GC_3) \simeq \mathfrak{grt}_1$
       (Theorem~\ref{thm:willwacher}; Definition~\ref{def:graph-complex}),
       acting on graph-integral weights $W_\Gamma$ (Construction~\ref{constr:graph-integral}).
 \item\label{S7} Stage-$1$ $\PhiFA_3(\cC)$ assembles local observables
       on holomorphic polydiscs via the $\FM_3$-model
       (Construction~\ref{constr:phifa3-via-fm};
       Theorem~\ref{thm:fm-to-disks}).
 \item\label{S8} The holomorphic Dolbeault refinement
       (Definition~\ref{def:dolbeault-fm}) upgrades the topological
       $E_3$-factorisation algebra to the holomorphic target of
       Costello--Li (Theorem~\ref{thm:cl-refinement}).
 \item\label{S9} Stage-$2$ $\SpCh_{\Sigma_2, C}$ degenerates the
       $\FM_3$-model to $\FM_1$ on $C$ by push-forward along a
       projection fibration (Construction~\ref{constr:fm3-to-fm1};
       Proposition~\ref{prop:e1-residual}).
 \item\label{S10} Dunn additivity $\FM_3 \simeq \FM_2 \otimes \FM_1$
       is a chain-level quasi-isomorphism (Theorem~\ref{thm:fm-dunn};
       Corollary~\ref{cor:compat-e3-obstr}), compatible with the
       target-chapter $E_3$-obstruction theorem.
\end{enumerate}
\end{survivor}

%==============================================================================
\section{Paths back to \texttt{cy\_categories.tex}}
%==============================================================================

The $\Chriss$--$\Ginzburg$-compliant inscription for
\texttt{chapters/theory/cy\_categories.tex} adds a new subsection to
Section~\texttt{cy-cat-phi-interface} recording:
\begin{itemize}
 \item a Definition block for the Fulton--MacPherson operad $\FM_n$
       on $\mathbb{R}^n$ with explicit boundary stratification;
 \item a Theorem block stating Kontsevich formality as a
       $\GRT_1(\mathbb{Q})$-torsor over $\Gerst_n \xrightarrow{\sim} C_n$;
 \item a Construction block exhibiting $\PhiFA_3(\cC)$ as an
       $\FM_3$-colimit of $\HH_\bullet(\cC)^{\otimes k}$-assemblies;
 \item a Proposition block recording the chain-level Dunn-additive
       decomposition $\FM_3 \simeq \FM_2 \otimes \FM_1$;
 \item a Remark block linking the $\FM_3$-content to the existing
       $\FM_k(C)$-reference in \texttt{rem:cy-cat-three-hochschild}
       (the chiral Gerstenhaber bracket from OPE residues on
       $\FM_k(C)$ is the $\SpCh$-specialisation of the topological
       Gerstenhaber bracket from $\FM_3$).
\end{itemize}

%==============================================================================
\section{Ghost theorems (unresolved intrinsic scope)}
%==============================================================================

\begin{remark}[Ghost theorems]
\label{rem:ghost-theorems}
Three ghost theorems remain at the boundary of the surviving core:
\begin{enumerate}
 \item[(G1)] \emph{Smooth-manifold Dunn additivity}: does the weak
             equivalence $\FM_{m+n} \simeq \FM_m \otimes_H \FM_n$ lift
             to a diffeomorphism of operads in smooth manifolds with
             corners? Present evidence suggests no (the smooth
             structures differ), but the precise obstruction lives in
             relative tangent complexes of moduli of operad models.
 \item[(G2)] \emph{$\GRT_1$-torsor transversality}: does every
             $\GRT_1(\mathbb{Q})$-element $\xi$ act on a fixed
             holomorphic-volume-trivialised $\PhiFA_3(\cC)$ output by
             a chain-level quasi-isomorphism that commutes with the
             Stage-$2$ $\SpCh_{\Sigma_2, C}$ specialisation on the
             nose, or only after a re-scaling of the specialisation
             datum? If the latter, the $\GRT_1$-action parametrises
             families of Stage-$2$ outputs rather than isomorphisms
             of a fixed output.
 \item[(G3)] \emph{Super-$\GRT_1$ extension}: in the Vol~III
             Etingof--Kazhdan quantisation of the K3 chiral bialgebra
             $\mathbf H_{\Delta_5}$ the relevant torsor is
             $\GRT_1^{\mathrm{super}}$, extending the bosonic
             $\GRT_1$ by a super-reflection \cite[AP-CY42,
             \texttt{notes/antipatterns\_catalogue.md}]{VolIII-APCY}.
             Does the chain-level $\FM_3$-model admit a super-refinement
             $\FM_3^{\mathrm{super}}$ whose $H^0(\GC_3^{\mathrm{super}})$
             is $\mathfrak{grt}_1^{\mathrm{super}}$?
\end{enumerate}
\end{remark}

\begin{thebibliography}{99}

\bibitem{FultonMacPherson1994}
W.\ Fulton and R.\ MacPherson.
\emph{A compactification of configuration spaces.}
Ann.\ of Math.\ (2) \textbf{139} (1994), no.\ 1, 183--225.

\bibitem{Kontsevich1999}
M.\ Kontsevich.
\emph{Operads and motives in deformation quantization.}
Lett.\ Math.\ Phys.\ \textbf{48} (1999), no.\ 1, 35--72.

\bibitem{Kontsevich2003}
M.\ Kontsevich.
\emph{Deformation quantization of Poisson manifolds.}
Lett.\ Math.\ Phys.\ \textbf{66} (2003), no.\ 3, 157--216.

\bibitem{KontsevichSoibelman2000-deformations}
M.\ Kontsevich and Y.\ Soibelman.
\emph{Deformations of algebras over operads and the Deligne conjecture.}
Conf.\ Mosh\'e Flato 1999, Vol.\ I (Dijon), 255--307, Math.\ Phys.\ Stud.\ 21, Kluwer, 2000.

\bibitem{BoardmanVogt1973}
J.\ M.\ Boardman and R.\ M.\ Vogt.
\emph{Homotopy invariant algebraic structures on topological spaces.}
Lecture Notes in Mathematics, Vol.\ 347. Springer, Berlin-Heidelberg, 1973.

\bibitem{BergerMoerdijk2006}
C.\ Berger and I.\ Moerdijk.
\emph{The Boardman--Vogt resolution of operads in monoidal model categories.}
Topology \textbf{45} (2006), no.\ 5, 807--849.

\bibitem{Salvatore2001}
P.\ Salvatore.
\emph{Configuration spaces with summable labels.}
In: Cohomological methods in homotopy theory (Bellaterra, 1998), 375--395, Progr.\ Math.\ 196, Birkh\"auser, 2001.

\bibitem{Sinha2004}
D.\ Sinha.
\emph{Manifold-theoretic compactifications of configuration spaces.}
Selecta Math.\ (N.S.) \textbf{10} (2004), no.\ 3, 391--428.

\bibitem{Sinha2006}
D.\ Sinha.
\emph{The topology of spaces of knots: cosimplicial models.}
Amer.\ J.\ Math.\ \textbf{131} (2009), no.\ 4, 945--980.

\bibitem{MarklShniderStasheff2002}
M.\ Markl, S.\ Shnider and J.\ Stasheff.
\emph{Operads in algebra, topology and physics.}
Mathematical Surveys and Monographs, 96. American Mathematical Society, Providence, RI, 2002.

\bibitem{Willwacher2015}
T.\ Willwacher.
\emph{M.\ Kontsevich's graph complex and the Grothendieck--Teichm\"uller Lie algebra.}
Invent.\ Math.\ \textbf{200} (2015), no.\ 3, 671--760.

\bibitem{Tamarkin2003}
D.\ E.\ Tamarkin.
\emph{Formality of chain operad of little discs.}
Lett.\ Math.\ Phys.\ \textbf{66} (2003), no.\ 1--2, 65--72.

\bibitem{HLTV2011}
R.\ Hardt, P.\ Lambrechts, V.\ Turchin and I.\ Voli\'c.
\emph{Real homotopy theory of semi-algebraic sets.}
Algebr.\ Geom.\ Topol.\ \textbf{11} (2011), no.\ 5, 2477--2545.

\bibitem{Dunn1988}
G.\ Dunn.
\emph{Tensor product of operads and iterated loop spaces.}
J.\ Pure Appl.\ Algebra \textbf{50} (1988), no.\ 3, 237--258.

\bibitem{LurieHA}
J.\ Lurie.
\emph{Higher Algebra.}
Preprint, 2017.

\bibitem{FrancisGaitsgory2012}
J.\ Francis and D.\ Gaitsgory.
\emph{Chiral Koszul duality.}
Selecta Math.\ (N.S.) \textbf{18} (2012), no.\ 1, 27--87.

\bibitem{FadellNeuwirth1962}
E.\ Fadell and L.\ Neuwirth.
\emph{Configuration spaces.}
Math.\ Scand.\ \textbf{10} (1962), 111--118.

\bibitem{Fresse2017}
B.\ Fresse.
\emph{Homotopy of operads and Grothendieck--Teichm\"uller groups.}
Mathematical Surveys and Monographs, 217. American Mathematical Society, Providence, RI, 2017.

\bibitem{CostelloGwilliam2017}
K.\ Costello and O.\ Gwilliam.
\emph{Factorization algebras in quantum field theory, Vol.\ 1.}
New Mathematical Monographs, 31. Cambridge University Press, 2017.

\bibitem{CostelloGwilliam2021}
K.\ Costello and O.\ Gwilliam.
\emph{Factorization algebras in quantum field theory, Vol.\ 2.}
New Mathematical Monographs, 41. Cambridge University Press, 2021.

\bibitem{CostelloLi2016}
K.\ Costello and S.\ Li.
\emph{Twisted supergravity and its quantization.}
arXiv:1606.00365 (2016).

\bibitem{CostelloLi2020}
K.\ Costello and S.\ Li.
\emph{Anomaly cancellation in the topological string.}
Adv.\ Theor.\ Math.\ Phys.\ \textbf{24} (2020), no.\ 7, 1723--1771.

\bibitem{CFG2026}
K.\ Costello, J.\ Francis and O.\ Gwilliam.
\emph{Factorization algebras for holomorphic Chern--Simons theories.}
Preprint, 2026.

\bibitem{GetzlerJones1994}
E.\ Getzler and J.\ D.\ S.\ Jones.
\emph{Operads, homotopy algebra, and iterated integrals for double loop spaces.}
Preprint, 1994. arXiv:hep-th/9403055.

\bibitem{Arnold1969}
V.\ I.\ Arnold.
\emph{The cohomology ring of the colored braid group.}
Mat.\ Zametki \textbf{5} (1969), 227--231.

\bibitem{Tamarkin2007}
D.\ Tamarkin.
\emph{What do dg-categories form?}
Compos.\ Math.\ \textbf{143} (2007), no.\ 5, 1335--1358.

\bibitem{VolIII}
\emph{Vol III: Calabi--Yau Quantum Groups.}
Manuscript in preparation. Chapter \texttt{chapters/theory/cy\_to\_chiral.tex}.

\bibitem{VolIII-APCY}
\emph{Vol III antipatterns catalogue.}
\texttt{notes/antipatterns\_catalogue.md}.

\end{thebibliography}

\end{document}
